\sect{AI-driven Robot Autonomy}{aidriven}

The evolution from conventional robot control to adaptive autonomy has been shaped by the increasing complexity and variability of real-world applications. Classical approaches based on deterministic models and predefined control laws, such as PID or trajectory-based programming, have proven insufficient for dynamic or unstructured environments \cite{isreal2025rise}. AI has therefore emerged as a key enabler of autonomy, allowing robotic systems to perceive, reason, and act with greater adaptability and independence \cite{machines13070615,pandy2025ai}. By introducing data-driven and learning-based components into the control architecture, AI-driven methods have enhanced a robot’s ability to interpret its surroundings, plan actions, and respond to uncertainty in real time.

Modern robotic autonomy builds on the integration of machine learning, probabilistic reasoning, and knowledge-based decision-making to achieve flexible task execution \cite{isreal2025rise,machines13070615}. Reinforcement learning and deep-learning-based planners have advanced navigation, manipulation, and coordination tasks in controlled domains, while hybrid systems combining symbolic reasoning with data-driven models have enabled hierarchical planning and context-aware operation \cite{pandy2025ai}. Despite these advances, such approaches remain limited in transparency and accessibility, as they typically require extensive data, complex training, and expert supervision. These constraints hinder their deployment in collaborative or human-centered environments where intuitive interaction and interpretability are essential.

Consequently, recent developments have shifted toward integrating NL based interfaces to bridge human intent and robotic execution. This transition marks the emergence of systems that combine high-level reasoning capabilities with linguistic understanding, enabling robots to follow unstructured instructions rather than predefined command sequences. This paradigm forms the foundation for NL Facilitated Human-Robot Cooperation, in which language models act as mediators between human objectives and autonomous system control.
